% Copy-ready replacement for the residual-attachment section of Version 2.
% Promote this text only after the literal attained-attachment identity and the
% q-only two-regime estimate have been validated on one integrated branch.

\section{Residual attachments}
\label{sec:residual-attachments-v2}

Section~\ref{sec:canonical-high-cells-v2} bounds the sum of the bare
high-skeleton weights.  It does not include the local rewards of the unexposed
cells or the binary cycle-space factor of the residual support.  We now bound
these remaining terms after conditioning on one canonical high skeleton.

Fix a feasible high skeleton with exposed block matching $M$.  Let $m_0$ be the
number of unexposed row stubs, equivalently the number of unexposed column
stubs.  Conditional on the exposed pairs, the remaining pairs form a uniform
bipartite matching with the prescribed residual degrees.

Let $\mathcal A(M)$ denote the conditional residual contribution, including the
local rewards of unexposed cells, the cycle-space factor, the residual
multiplicity cap, and the condition that no exposed cell receives another
pair.  The normalized second moment has the exact decomposition
\begin{equation}
  \frac{\mathbb E Z^2}{(\mathbb E Z)^2}
  =
  \sum_M w_{\mathrm{hi}}(M)\,\mathcal A(M),
  \label{eq:exact-attachment-decomposition-v2}
\end{equation}
where the sum runs over the finite family of feasible canonical high
skeletons.  Thus it remains to bound $\mathcal A(M)$ uniformly.

\subsection{Restriction outside the exposed matching}

The binary cycle-space identity writes the cycle contribution as a sum over
even residual edge sets.  The following finite lemma replaces a decomposition
into simple cycles.

\begin{lemma}[Restriction-product bound]
\label{lem:restriction-product-v2}
Let $E$ be a finite set, let $\mathcal C$ be a finite family of subsets of
$E$, and let $I\subseteq E$.  Suppose that the map
\[
  A\longmapsto A\setminus I
\]
is injective on $\mathcal C$.  For nonnegative activities $(q_e)_{e\in E}$,
\[
  \sum_{A\in\mathcal C}
    \prod_{e\in A\setminus I}q_e
  \le
  \prod_{e\in E\setminus I}(1+q_e).
\]
\end{lemma}

\begin{proof}
The sets $A\setminus I$, with $A\in\mathcal C$, form a subfamily of the power
set of $E\setminus I$ and occur without repetition.  Enlarging the sum to the
entire power set gives
\[
  \sum_{B\subseteq E\setminus I}\prod_{e\in B}q_e
  =
  \prod_{e\in E\setminus I}(1+q_e).
\]
\end{proof}

Apply the lemma to the family of even edge sets and to the exposed matching
$I=M$.  The required injectivity is immediate.  If two even edge sets have the
same restriction outside $M$, then their symmetric difference is an even
subset of the matching $M$.  A nonempty subset of a matching has a vertex of
degree one, so no such subset is even.  Therefore the symmetric difference is
empty.

It follows that
\begin{equation}
  \sum_{F\ \mathrm{even}}
    \prod_{e\in F\setminus M}q_e
  \le
  \prod_{e\notin M}(1+q_e).
  \label{eq:even-family-product-v2}
\end{equation}
This estimate is exact at the level of finite supports and does not require a
choice of cycle decomposition.

\subsection{One activity controls both residual factors}

For a residual cell $e$, let $\lambda_e^{\mathrm{loc}}$ be the contribution of
local multiplicity increments and let $q_e$ be the total activity used in the
even-subgraph expansion.  The definitions give the pointwise bound
\[
  0\le \lambda_e^{\mathrm{loc}}\le q_e.
\]
Consequently
\[
  \prod_e(1+\lambda_e^{\mathrm{loc}})
  \prod_{e\notin M}(1+q_e)
  \le
  \exp\!\left(2\sum_e q_e\right).
\]
Thus the local rewards and the cycle-space factor are controlled by one
quantity, namely the total residual activity.

The residual estimate has two regimes.  Let $U$ be the phase cap.

\begin{proposition}[Uniform residual attachment bound]
\label{prop:q-only-attachment-v2}
There is an absolute constant $C$ such that the following holds for every
feasible canonical high skeleton.

If
\[
  2^U\le m_0^3,
\]
then
\[
  \mathcal A(M)\le \exp(CU^2).
\]
If $2^U>m_0^3$, then
\[
  m_0<2^{\lceil U/3\rceil}
\]
and
\[
  \mathcal A(M)\le 2^{Um_0/2}.
\]
\end{proposition}

\begin{proof}
In the first regime, the residual degree bounds and the definition of the
activities give
\[
  \sum_e q_e\le C_1U^2.
\]
The preceding product estimate yields
$\mathcal A(M)\le\exp(2C_1U^2)$.

In the second regime, the inequality $2^U>m_0^3$ implies
$m_0<2^{\lceil U/3\rceil}$.  We discard the residual restrictions and use the
deterministic bounds
\[
  \sum_e\binom{r_e}{2}
  \le \frac{U-1}{2}m_0,
  \qquad
  \beta(H_{\mathrm{res}}\cup M)
  \le \frac{m_0}{2}.
\]
The product of the local and cycle-space factors is therefore at most
$2^{Um_0/2}$.
\end{proof}

The two estimates are uniform in the canonical high skeleton.  For the exact
midpoint profile, they imply the following aggregate form.

\begin{proposition}[Attained attachment sum]
\label{prop:attained-attachment-sum-v2}
There is a deterministic sequence $\varepsilon_n\to0$ such that
\begin{equation}
  \operatorname{AttachmentSum}_n
  \le
  \operatorname{BareSkeletonSum}_n
  \exp\!\left\{
    \varepsilon_n\frac{n}{(\log n)^4}
  \right\}.
  \label{eq:attained-attachment-sum-v2}
\end{equation}
\end{proposition}

\begin{proof}
Use the first bound in Proposition~\ref{prop:q-only-attachment-v2} when
$2^U\le m_0^3$.  Since $U=O(\log n)$, its logarithm is $O((\log n)^2)$.
In the complementary regime, the inequality
$m_0<2^{\lceil U/3\rceil}$ and the phase relation imply
\[
  Um_0=o\!\left(\frac{n}{(\log n)^4}\right).
\]
Taking the larger of the two deterministic errors gives a sequence
$\varepsilon_n\to0$ that is independent of the skeleton.  Factoring this
uniform bound from the exact decomposition
\eqref{eq:exact-attachment-decomposition-v2} proves
\eqref{eq:attained-attachment-sum-v2}.
\end{proof}

Combining Proposition~\ref{prop:attained-attachment-sum-v2} with
\eqref{eq:bare-skeleton-final-ams-v2} yields
\[
  \log\frac{\mathbb E Z^2}{(\mathbb E Z)^2}
  =
  o\!\left(\frac{n}{(\log n)^4}\right).
\]
This is the normalized second-moment estimate used in the amplification
argument.
